% ======================================================================
\section{Introduction}\label{sec:intro}
% ======================================================================

\subsection{The problem and the two algorithms}

We want to solve the unconstrained minimization problem
\begin{equation}\label{eq:problem}
  \min_{x\in\R^d} f(x),
\end{equation}
where $f:\R^d\to\R$ is differentiable, but its gradient $\nabla f$ is not available
exactly. Instead, at any point $x$ we can ask for a \emph{stochastic gradient}
$g(x;\xi)$: a random vector, depending on a random sample $\xi$, which is correct
\emph{on average}, $\E_\xi[g(x;\xi)]=\nabla f(x)$. This is the situation in machine
learning, where $f(x)=\frac1n\sum_{i=1}^n f_i(x)$ is the average loss over $n$ training
examples and $g(x;\xi)=\nabla f_\xi(x)$ is the gradient of the loss on one randomly chosen
example (or on a random mini-batch).

\paragraph{SGD.} Stochastic gradient descent starts at a point $x_0$ and repeats
\begin{equation}\label{eq:sgd}
  x_{k+1}=x_k-\alpha\, g_k,\qquad g_k:=g(x_k;\xi_k),\qquad k=0,1,2,\dots
\end{equation}
where $\alpha>0$ is the \emph{step size} (learning rate) and $\xi_0,\xi_1,\dots$ are
independent samples.

\paragraph{SGDM.} SGD with momentum---also called the \emph{stochastic heavy-ball
method}, after Polyak's heavy-ball method~\cite{polyak1964}---adds a fraction
$\beta\in[0,1)$ of the previous displacement:
\begin{equation}\label{eq:sgdm-intro}
  x_{k+1}=x_k-\alpha\, g_k+\beta\,(x_k-x_{k-1}),\qquad x_{-1}:=x_0 .
\end{equation}
Physically, $x_k$ is the position of a heavy ball rolling down the graph of $f$: the term
$\beta(x_k-x_{k-1})$ is its inertia (it keeps moving in the direction it was already
moving), and $1-\beta$ plays the role of friction. With $\beta=0$ we recover SGD. In deep
learning, \eqref{eq:sgdm-intro}---usually implemented in an equivalent ``momentum buffer''
form (\cref{sec:forms})---with $\beta=0.9$ is a standard default~\cite{sutskever2013}.

\subsection{What ``convergence'' means here}

For a nonconvex $f$ there is no efficient method that is guaranteed to find a global
minimizer, so the standard goal is an approximate \emph{stationary point}, a point with a
small gradient. Since the iterates are random, guarantees are stated in expectation. A
typical statement has the form
\[
  \min_{0\le k<K}\E\norm{\nabla f(x_k)}^2\;\le\;(\text{a quantity that tends to $0$ as
  $K\to\infty$}),
\]
and the speed at which the right-hand side decreases in the number of iterations $K$ is
the \emph{rate}. For special classes of functions (Polyak--{\L}ojasiewicz or convex) we
will also bound the optimality gap $\E[f(x_k)]-\fs$, where $\fs=\inf f$.

\subsection{The difficulty, and the idea in one line}

For SGD, every textbook proof~\cite{bottou2018,ghadimilan2013} rests on one inequality,
the \emph{expected decrease}:
\begin{equation}\label{eq:intro-sgd-decrease}
  \E\bigl[f(x_{k+1})\,\big|\,\text{history}\bigr]\;\le\;f(x_k)
  -\frac{\alpha}{2}\norm{\nabla f(x_k)}^2+\frac{L\alpha^2}{2}\sigma^2
  \qquad(\alpha\le 1/L),
\end{equation}
where $L$ is the smoothness constant of $f$ and $\sigma^2$ bounds the variance of the
stochastic gradients (\cref{sec:sgd}). In words: on average, one step decreases $f$ by an
amount proportional to the squared gradient, minus a price for the noise that is
\emph{quadratic} in the step size. Everything else---rates, step-size rules---follows by
summing \eqref{eq:intro-sgd-decrease} over $k$.

For SGDM, inequality \eqref{eq:intro-sgd-decrease} is \emph{false}, even without noise:
momentum can carry the iterate uphill (in \cref{ex:overshoot} the method sits exactly at
the minimizer and is pushed away by its own inertia). It holds only at steps where the
momentum still points downhill (\cref{prop:f-step}), and the algorithm cannot control
when that is the case. The function value is the wrong yardstick. The idea of this report is to measure progress with a \emph{potential}
\[
  \Phi_k=f(z_k)-\fs+\frac{L\beta^2}{2(1-\beta)^2}\norm{x_k-x_{k-1}}^2,
  \qquad z_k=x_k+\frac{\beta}{1-\beta}\,(x_k-x_{k-1}),
\]
where $z_k$ is an extrapolated point---the place where the ball would come to rest if all
forces were switched off (\cref{sec:extrapolated})---and the second term is a kinetic
energy. We prove (\cref{thm:master}) that $\Phi_k$ obeys the exact analogue of
\eqref{eq:intro-sgd-decrease}:
\begin{tcolorbox}[keybox]
\begin{equation}\label{eq:intro-master}
  \E\bigl[\Phi_{k+1}\,\big|\,\text{history}\bigr]\;\le\;\Phi_k
  -\frac{\gamma}{4}\norm{\nabla f(x_k)}^2+L\gamma^2\sigma^2,
  \qquad \gamma:=\frac{\alpha}{1-\beta},
\end{equation}
\centering whenever the step size satisfies $\displaystyle\alpha\le\frac{(1-\beta)^2}{2L}$.
\end{tcolorbox}
\noindent Once \eqref{eq:intro-master} is available, all the standard SGD conclusions
transfer to SGDM by the same telescoping arguments, with the \emph{effective step size}
$\gamma=\alpha/(1-\beta)$ playing the role of $\alpha$.

\subsection{Summary of results}

\Cref{tab:summary} lists what we prove. Throughout, $\Dlt:=f(x_0)-\fs$, $\gamma=\alpha/(1-\beta)$,
and ``nonconvex'' means only that $f$ is $L$-smooth and bounded below.

\begin{table}[ht]
\centering\small
\renewcommand{\arraystretch}{1.25}
\begin{tabularx}{\textwidth}{@{}l l >{\raggedright\arraybackslash}X l@{}}
\toprule
Setting & Step-size condition & Guarantee & Where\\
\midrule
Nonconvex & $\alpha\le\frac{(1-\beta)^2}{2L}$ &
  $\frac1K\sum_{k<K}\E\norm{\nabla f(x_k)}^2\le\frac{4\Dlt}{\gamma K}+4L\gamma\sigma^2$ &
  \cref{thm:nonconvex}\\
Nonconvex, tuned $\gamma$ & $\gamma=\min\bigl\{\frac{1-\beta}{2L},\frac1\sigma\sqrt{\frac{\Dlt}{LK}}\bigr\}$ &
  $\frac1K\sum_{k<K}\E\norm{\nabla f(x_k)}^2\le\frac{8L\Dlt}{(1-\beta)K}+8\sigma\sqrt{\frac{L\Dlt}{K}}$ &
  \cref{cor:rate}\\
Nonconvex, decaying & $\gamma_k\le\frac{1-\beta}{2L}$, $\sum\gamma_k=\infty$, $\sum\gamma_k^2<\infty$ &
  $\liminf_k\norm{\nabla f(x_k)}=0$ almost surely &
  \cref{thm:diminishing}\\
No noise ($\sigma=0$) & $\alpha=\frac{(1-\beta)^2}{2L}$ &
  $\Phi_k$ non-increasing; $\min_{k<K}\norm{\nabla f(x_k)}^2\le\frac{8L\Dlt}{(1-\beta)K}$ &
  \cref{cor:deterministic}\\
No noise, smooth & $\alpha<\frac{2(1-\beta)}{L}$ &
  energy decreases; $\nabla f(x_k)\to0$ &
  \cref{prop:energy}\\
PL inequality & $\alpha\le\frac{(1-\beta)^2}{2L}$ &
  $\E f(x_k)-\fs\le2\bigl(1-\frac{\gamma\mu}{4}\bigr)^k\Dlt+\frac{8L\gamma\sigma^2}{\mu}$ &
  \cref{thm:pl}\\
PL, decaying & $\gamma_k=\frac{8}{\mu(k+a)}$ &
  $\E f(x_k)-\fs=O(1/k)$ &
  \cref{cor:pl-decay}\\
Convex & $\alpha\le\frac{1-\beta}{2L}$ &
  $\E f(\bar x_K)-\fs\le\frac{\norm{x_0-x^\star}^2}{\gamma K}+\frac{2\beta\Dlt}{(1-\beta)K}+\gamma\sigma^2$ &
  \cref{thm:convex}\\
\bottomrule
\end{tabularx}
\caption{Summary of the guarantees proved in this report (all for SGDM in heavy-ball form
with momentum $\beta\in[0,1)$; for $\beta=0$ they reduce to the corresponding SGD results
up to constants). The question ``under what circumstances?'' is discussed in
\cref{sec:circumstances}: there we show that the noise floor $\propto\gamma\sigma^2$ is
unavoidable with a constant step (\cref{prop:noise-floor}), that too large a step size
makes the method diverge or cycle (\cref{prop:quadratic-stability},
\cref{ex:lrp}), and how the admissible step sizes shrink as $\beta\to1$
(\cref{fig:regions}).}
\label{tab:summary}
\end{table}

\subsection{How to read this report}

\paragraph{Prerequisites.} Multivariable calculus (gradients, the chain rule, the
fundamental theorem of calculus), linear algebra (inner products, norms, eigenvalues of
$2\times2$ matrices) and probability (expectation, independence, and conditional
expectation at the level of ``average over what is not yet known''). \Cref{sec:prelim}
collects every tool we use, with proofs; later proofs cite these facts explicitly.

\paragraph{Road map.} \Cref{sec:sgd} reviews the SGD analysis, which serves as the
template. \Cref{sec:momentum} introduces SGDM, explains why the function value is the
wrong yardstick, and introduces the extrapolated point $z_k$. \Cref{sec:master}, the heart
of the report, proves the expected decrease inequality \eqref{eq:intro-master}.
\Cref{sec:convergence} turns it into convergence theorems. \Cref{sec:circumstances}
answers the question ``under what circumstances does SGDM converge?'', including exact
computations that show which parts of the conditions cannot be removed.
\Cref{sec:experiments} illustrates everything numerically; the code that produces every
figure is in the \texttt{experiments/} folder of the repository.

\paragraph{Relation to the literature.} The extrapolated-point (``virtual sequence'')
technique and potentials of this type are classical in analyses of momentum
methods~\cite{ghadimi2015,yang2016,yu2019,liu2020,sebbouh2021}. The aim here is a
self-contained, fully explicit version---one potential, one key inequality, and explicit
constants---written so that each line can be checked by hand. Constants are chosen for
readability, not optimized.
